\documentclass[10pt,a4paper]{article}
\usepackage{onepager}

\begin{document}
\small

\ophead{Perception Pressure and Resonance}
{Decomposable, Trust-Weighted Measures of Media Narrative from a Multi-Perspective News-Aggregation System}
{\code{prism_preprint.tex} $\cdot$ 7pp $\cdot$ Jun 2026 $\cdot$ system lives at \code{~/Prism}
 (separate repo) $\cdot$ lead preprint for Leippold (UZH) $\cdot$ methods/system paper, no empirics}

\begin{multicols}{2}

\opsec{Thesis}
Two quantitative measures of media narrative computed continuously from a live multi-source news stream,
plus the system that produces them. Not another sentiment index: the design goal is a measure that is
\emph{signed}, \emph{magnitude-meaningful}, \emph{manipulation-resistant}, and \emph{fully auditable}.

\opsec{Resonance — per-story media impact}
\[
  \mathrm{Res}(c,t)=\underbrace{\mathrm{breadth}(c)}_{\text{diversity multiplier}}\cdot
   \sum_{m\in c}\mathrm{trust}(s_m)\,\mathrm{eng}(m)\,\mathrm{dec}(a_m)
\]
with
$\mathrm{breadth}(c)=\log_2(1+n_{\text{sources}})$,
$\mathrm{eng}(m)=1+\log_2(1+\mathrm{re}_m/\text{median})$,
$\mathrm{dec}(a)=e^{-\lambda a}$, $\lambda=\ln2/h$.

\opsec{Perception Pressure — per-concept longitudinal signal}
\begin{align*}
 \text{salience:}\ & A(K,t)=\textstyle\sum_c \mathrm{breadth}(c)\,\mathrm{dec}(a_c)\ \ge 0\\
 \text{valence:}\ & V(K,t)\in[-1,+1]\ \text{(credibility-weighted)}\\
 \text{perception:}\ & \boxed{\,P(K,t)=A(K,t)\cdot V(K,t)\,}
\end{align*}
momentum $\dot P(K,t)=P(K,t)-P(K,t-1)$.

\opsec{Three properties, novel \emph{in combination}}
\begin{opnums}
\item \textbf{Multiplicative coupling} of attention and directional sentiment $\Rightarrow$ the signal is
      \emph{signed} and its magnitude means something. Volume-only proxies have no sign; sentiment-only
      indices have no weight.
\item \textbf{Diversity as a multiplier, not a sum} $\Rightarrow$ a single outlet cannot inflate the
      measure. Multi-perspective by construction, not by editorial policy.
\item \textbf{Every component persisted independently} $\Rightarrow$ fully decomposable and auditable.
      You can always answer ``why did $P$ move'' — it is not a black box.
\end{opnums}

\opsec{Trust is earned, not assigned}
\[
  \mathrm{trust}=0.1+0.4\cdot\frac{\#\,\text{validated}}{\#\,\text{total}}
\]
via an automated cross-corroboration lifecycle. \textbf{No hand-assigned credibility table} — which is
where comparable indices smuggle in the researcher's politics, and the first thing a referee will
probe.

\opsec{The system}
Five-agent pipeline — \emph{discover $\to$ analyze $\to$ track $\to$ personalize $\to$ deliver} —
coordinated by a database state machine. Prism is the upstream producer of the narrative/hype series NAT
ingests.

\opsec{Positioning}
Against volume-only attention proxies (Google Trends, article counts) and direction-only sentiment
indices. The claim is not ``better sentiment'' but \emph{a different object}: a signed pressure with
interpretable units and an audit trail.

\opsec{Downstream application (specified, not claimed)}
Treat $P(K,t)$ as an \textbf{exogenous narrative signal} whose incremental predictive value for financial
returns is measured under a leakage-controlled information-theoretic protocol.
\caution{The paper is explicit that this is a methods-and-system paper: the instrument is fully specified
and parameterized for ablation; \textbf{empirical validation is the next step, not a claim made here}.
Neither Prism metric has logged results. Do not let a reader infer otherwise.}

\opsec{Why it is honest to publish results-free}
Same posture as \#3 and \#5: specify the instrument, expose the constraint, then measure. A referee can
attack the construction on its merits without having to disentangle it from a performance claim — and
there is nothing to retract when the empirics land.

\opsec{The five agents}
Coordinated not by a message queue but by a \textbf{database state machine}: \textsc{raw} $\to$
\textsc{analyzed} $\to$ \textsc{delivered}.
\begin{opfacts}
Discovery & pull + cluster articles by event; grow the source registry (RSS ingest; lexical dedup — Jaccard, TF--IDF, entity overlap) \\
Analysis & summarize a cluster; extract per-source perspective, sentiment, bias, claims (LLM structured extraction + quality score) \\
Tracking & scan clusters for tracked concepts; compute the measures \\
Personalize & per-reader selection \\
Deliver & emit \\
\end{opfacts}
Stated editorial principle: make bias \emph{visible} rather than hidden — no original reporting, every
claim attributed, multiple framings side by side. The quantitative measures are instrumentation of that
principle, not a bolt-on.

\opsec{Disclosed threats to validity}
\caution{Sentiment/bias labels come from an \textbf{LLM, not gold-standard human annotation} — the
structure of that measurement error is itself an open research question, and it is the first thing a
referee will press. Clustering is \textbf{lexical, not embedding-based} (token overlap, TF--IDF, regex
entity extraction), so events can split or merge. The engagement term \textbf{defaults to neutral} when
reaction counts are missing, so it is often inert in practice. Storage is single-writer, bounding
throughput. None of these undermine the \emph{definitions}; all of them bear on the fidelity of the
computed series and must be controlled in any empirical study.}

\opsec{Siblings}
\#5 (\emph{From Narrative to Alpha}) is the direct sequel: it takes $P(K,t)$ as given and argues, at the
level of methodology, what an exogenous signal contributes to a microstructure book — timescale
separation, effective breadth, partial evasion of the adverse-selection barrier. Send \#4 then \#5 to
Leippold; \#4 alone reads as a systems paper.

\end{multicols}

\opfoot{One-pager for personal use. The distinguishing feature to lead with: \textbf{diversity enters as
a multiplier}, so the measure is structurally resistant to single-outlet manipulation — and trust is
earned by corroboration rather than assigned by hand.}

\end{document}
